% =====================================================================
% Section 1. Introduction  (structure as Battisti & Boscheri 2025, Sec. 1)
% =====================================================================
\section{Introduction}
\label{sec:intro}

Compressible liquid--gas and liquid--vapour flows occur in cavitating hydraulic machinery and fuel injectors,
in underwater explosions, in the rapid depressurisation of pressurised water in nuclear reactor safety, and in
the shock loading of condensed materials and mixtures \cite{SaurelPantano2018,SaurelPetitpasAbgrall2008,Bartak1990}.
Two features make them difficult. First, the Mach number varies by orders of magnitude within one flow: a
liquid moving at a few metres per second is in the incompressible regime, while the gas, a vapour pocket or a
detonation product in the same domain may be transonic or carry strong shocks. Second, the phases exchange
momentum, heat and mass on time scales that range from much slower to much faster than the flow itself, so that
the same computation contains frozen, finite-rate and equilibrium regions, and phase transition through
metastable states.

Diffuse-interface models describe such flows with one set of equations in the whole domain, pure fluids,
interfaces and mixtures alike \cite{SaurelAbgrall1999a,AbgrallKarni2001,SaurelPantano2018}. They derive from the
seven-equation model of Baer and Nunziato \cite{BaerNunziato1986}, in which each phase has its own velocity and
pressure. Relaxing the velocities gives the single-velocity six-equation model with two pressures; relaxing also
the pressures gives the five-equation model of Kapila et al.\ \cite{Kapila2001,MurroneGuillard2005,AllaireClercKokh2002},
and further thermal and chemical relaxation give the four-equation $pT$ model and the homogeneous equilibrium model
\cite{FlattenLund2011,Lund2012,Pelanti2022}. The five-equation model is the natural limit for interfaces and mixtures,
but its volume-fraction equation carries a non-conservative compressibility term that makes volume fractions
lose positivity in strong expansions and compressions, and its mixture sound speed, the Wood speed, is not
monotone in the volume fraction. Saurel, Petitpas and Berry \cite{SPB2009} therefore proposed to solve the
six-equation model with stiff pressure relaxation: its frozen sound speed is monotone, the volume-fraction
equation contains only the pressure work, and the mixture total energy, which is redundant for smooth flows, is
added as a seventh equation (``6+1'') to restore the conservative mixture energy at shocks. Pelanti and Shyue
\cite{PelantiShyue2014} made the scheme mixture-energy-consistent and added heat and mass transfer through
thermal and chemical relaxation, following Saurel et al.\ \cite{SaurelPetitpasAbgrall2008} and Zein et al.\
\cite{Zein2010}; Pelanti \cite{Pelanti2022} extended these procedures to arbitrary relaxation rates and to the
Noble--Abel stiffened-gas \cite{LeMetayerSaurel2016} and Mie--Gr\"uneisen equations of state. The six-equation
relaxation approach is now a standard tool for cavitation, flashing and evaporation fronts
\cite{PelantiShyue2019,Chiapolino2017,DeLorenzo2018,DeLorenzo2019,Schmidmayer2020,SirianniRe2026}. A known difficulty
is the behaviour at shocks in mixtures: the six-equation model is non-conservative, its shock relations
\cite{SaurelGavrilyukRenaud2003} are not those of the scheme's numerical shock layer, and the energy deposited
in the layer is shared between the phases in a way that depends on the scheme. Saurel et al.\ \cite{SPB2009}
showed that the computed post-shock state then deviates from the five-equation solution for strong shocks, and
corrected it with an artificial heat exchange in the shock layer whose coefficient must be calibrated for each
mixture \cite{Petitpas2007}.

All the numerical methods above are explicit Godunov-type schemes with HLLC, Roe or wave-propagation
solvers for the hyperbolic part, followed by relaxation steps
\cite{SPB2009,PelantiShyue2014,Pelanti2022,DeLorenzo2018,SirianniRe2026}. Their time step is limited by the fastest
acoustic speed. In a single-velocity mixture this is the frozen sound speed, which a small amount of liquid
raises to nearly the liquid value, so that a flow with velocities of a few metres per second is advanced with the
time step of a wave that moves at more than a kilometre per second. In addition, the numerical dissipation of upwind
schemes scales with the sound speed and destroys the low-Mach asymptotics of the pressure
\cite{GuillardViozat1999,GuillardMurrone2004,Dellacherie2010}. Preconditioning restores the correct scaling of the
dissipation \cite{Turkel1987,Pelanti2017} and, combined with implicit time integration, removes the time step
restriction \cite{LeMartelot2013}, but at the price of large nonlinear systems coupling all variables, whose
solution is costly and whose convergence is difficult to control.

Treating implicitly only the terms responsible for the stiffness, while keeping the nonlinear transport
explicit, avoids both difficulties. Pressure-based semi-implicit methods go back to Harlow and Welch
\cite{HarlowWelch1965} and Harlow and Amsden \cite{HarlowAmsden1971}, with fundamental contributions by Casulli and
Greenspan \cite{CasulliGreenspan1984}, Klein \cite{Klein1995}, Munz et al.\ \cite{Munz2003} and Park and Munz
\cite{ParkMunz2005}. Conservative, shock-capturing versions are obtained by splitting the flux into a convective
part with eigenvalues given by the flow velocity and a pressure part carrying the acoustic waves
\cite{ZhaBilgen1993,LiouSteffen1993,ToroVazquez2012}. Dumbser and Casulli \cite{DumbserCasulli2016} combined the
Toro--V\'azquez-Cend\'on splitting with an implicit pressure sub-system for general equations of state; the
resulting system is only mildly nonlinear, the nonlinearity being confined to the diagonal through the equation of
state, and is solved by the nested Newton method of Casulli and Zanolli \cite{CasulliZanolli2012}. Higher-order and
all-Mach extensions followed in finite volume and discontinuous Galerkin frameworks
\cite{TavelliDumbser2017,Boscheri2020,BoscheriPareschi2021,BoscheriTavelli2022,Boscarino2022,Orlando2022}. Time
integration with second and higher order is provided by implicit--explicit (IMEX) Runge--Kutta schemes
\cite{Ascher1997,PareschiRusso2005}, whose stiffly accurate, L-stable members preserve relaxation limits
\cite{PareschiRusso2005,BoscarinoRusso2009,BoscarinoPareschiRusso2017}, and by their semi-implicit and linearly
implicit variants \cite{BoscarinoFilbetRusso2016,Avgerinos2019}. A scheme that reduces, as the Mach number tends to
zero, to a consistent discretisation of the incompressible limit with a time step independent of the Mach number
is asymptotic preserving \cite{Jin1999,KlainermanMajda1981}; such all-speed schemes have been analysed for the Euler
equations in \cite{DegondTang2011,HaackJinLiu2012,Cordier2012,Noelle2014,Dimarco2018,Boscheri2020}. The same
construction carries over to systems with more wave families: in ideal magnetohydrodynamics the Alfv\'en and
magnetosonic waves have been removed successively from the explicit step
\cite{Dumbser2019,Fambri2021,BoscheriThomann2024,Dematte2024}, and a pressure-based formulation with a general
equation of state and pressure positivity enforced inside the implicit solve was given in \cite{Chen2026}.

For two-phase flows, all-Mach methods are still few. Re and Abgrall \cite{ReAbgrall2022} proposed a pressure-based
formulation of a Baer--Nunziato type model for weakly compressible flows with general equations of state, and Rana
and Singh \cite{RanaSingh2024} a pressure-based solver that switches between compressible and incompressible
behaviour; Varsakelis and Papalexandris \cite{VarsakelisPapalexandris2014} prescribe the compressibility of each
phase a priori. Demou et al.\ \cite{Demou2022} developed a pressure-based diffuse-interface method for low-Mach flows
with mass transfer. Flux splittings separating convection and acoustics have been derived for the Baer--Nunziato
\cite{TokarevaToro2016} and Kapila \cite{tenEikelder2017} models, and IMEX schemes have been constructed for a
two-fluid single-temperature model \cite{Lukacova2024} and for the full Baer--Nunziato model
\cite{MalusaAlaia2024}. Battisti and Boscheri \cite{BattistiBoscheri2025} designed a linearly implicit, asymptotic
preserving and well-balanced scheme for the seven-equation model, with explicit Rusanov fluxes for the convective
terms, central implicit discretisation of the pressure terms and implicit stiff relaxation sources. Closest to
the present work, Schropff, Petitpas and Daniel \cite{Schropff2026} derived a Zha--Bilgen type splitting of the
single-velocity two-pressure model and a semi-implicit scheme in which the pressure sub-system is solved by Picard
iterations and the phase internal energies are corrected to conserve total energy; their scheme is restricted to
one space dimension and instantaneous pressure relaxation, and the dissipation of its explicit step is a multiple of
the material speed with a coefficient chosen for each test. To our knowledge, no IMEX scheme exists for the
six-equation model with mechanical, thermal and chemical relaxation, which is the model used in practice for
cavitating and evaporating flows.

In this work we construct a second-order IMEX finite volume scheme for the six-equation model with total energy
and the three relaxations, for stiffened-gas, Noble--Abel stiffened-gas and Mie--Gr\"uneisen equations of state, in
one and two space dimensions. The time step is set by a CFL condition on the material speed only, whatever the Mach
number and the relaxation rates. The main ingredients are the following.
\begin{itemize}
\item The fluxes are split following Toro and V\'azquez-Cend\'on \cite{ToroVazquez2012}: the explicit part
transports all phase quantities with the old velocity and a Rusanov flux whose dissipation speed is the material
speed, and the implicit part carries the pressure, its work on each phase and an incremental flux of the
transported quantities. The volume-fraction equation is written in volume-transport form, so that its
non-conservative product becomes a conservative volume flux plus the volume strain.
\item The implicit pressure work and the three relaxation processes enter a single local closure, which for
each set of rates, finite or infinite, reduces to one scalar equation for the cell pressure. The stage then
reduces to a nonlinear system for the cell pressures, mildly nonlinear in the sense of
\cite{DumbserCasulli2016,CasulliZanolli2012}, with a three-point (five-point in 2D) Jacobian, solved by a Newton
method whose iteration count is observed to be independent of the Mach number.
\item The difference between the conservative and the non-conservative mixture internal energy, which at a
shock is the dissipation of the scheme, is distributed between the phases according to the shock relations of the
model \cite{SaurelGavrilyukRenaud2003}: the physical part in proportion to the Hugoniot--isentrope energy gap of
each phase, the discretisation error in the unique linear proportion that keeps both phases admissible. No
calibrated heat exchange is needed.
\item The acoustic dissipation that the central implicit operator lacks is supplied at the faces by the
acoustic Riemann solution, weighted by the deficit of the implicit viscosity with respect to the upwind viscosity,
and limited so that the scheme is second order in smooth regions and local-extremum-diminishing.
\item Second order in time is obtained with the stiffly accurate, L-stable LSDIRK2 scheme
\cite{BoscarinoFilbetRusso2016}; the start-up steps use one backward-Euler stage.
\end{itemize}
We prove that, under non-degeneracy assumptions, the scheme converges to the corresponding equilibrium scheme as each
relaxation rate tends to infinity and reduces to the scheme without that relaxation as the rate vanishes; that it
preserves moving contacts (with equal phase temperatures when thermal relaxation is active) and the temperature of each
phase under transport; and that, for stiffened and Noble--Abel stiffened gases, the transported phase states of a
backward-Euler stage remain admissible under a face condition independent of the phase contents. We show formally
that the scheme is asymptotic preserving in the low Mach number limit of the mixture. The scheme is validated on
twenty-two problems in one and two space dimensions (counting the single- and two-phase Taylor--Green vortices
separately): the convergence and low-Mach tests of \cite{BattistiBoscheri2025}, single-phase Riemann problems at low and
high Mach numbers and a low-Mach interface advection from \cite{Schropff2026}, and the interface, shock,
cavitation, evaporation and depressurisation problems of \cite{SPB2009,PelantiShyue2014,Pelanti2022}.

The rest of the paper is organised as follows. Section~\ref{sec:governing} presents the mathematical model, the
equations of state and the relaxed limit models. Section~\ref{sec:method} describes the numerical scheme and proves
its properties. Section~\ref{sec:results} presents the numerical results in one and two space dimensions, and
conclusions are drawn in Section~5.
